% SPDX-License-Identifier: Apache-2.0
\section{Blinky}
\label{sec:x-blinky}

One register, one wire, one \code{when!}, and a build that says how
fast the clock is and how fast to blink. The period is an associated
constant computed from the two rates, so the board build and the
simulation build differ in two numbers. The LED is derived from the
count in the same iteration that reads it, so there is one register
where a Verilog blinky often holds two. There is no top: \code{main}
splits the wire, the blinky takes the driving end by value as its
output, and the testbench keeps the reading end and prints it once per
cycle.

The blinky is also the first unit with state to be lowered.
\code{\#[lower]} on its \code{Unit} impl reads the \code{run} loop,
one \code{rising().await}, a register read, a \code{when!} and a drive,
and writes \code{Blinky::verilog(name)} beside it. The registers and
their widths come from the struct, through the same \code{Trace} derive
that names signals for a waveform, and the constants of the
configuration are evaluated when \code{verilog} runs, so the two
builds lower to two modules from one unit: the simulation build
compares the count with 15 and 8, the board build with 99999999 and
50000000. Both modules are at the end of Listing~\ref{lst:x-blinky-out}.
The comparison had to be written \code{lt(n, U::from(C::HALF))} rather
than on the raw integers for the macro to read it, which is the first
place the lowering pushed back on the source.

\lstinputlisting[caption={\code{ex\_blinky.rs}. Run with \code{bazel run //lib/examples:ex\_blinky}.},label={lst:x-blinky}]{lib/examples/ex_blinky.rs}

\lstinputlisting[style=ascii,caption={What \code{ex\_blinky} printed when the build ran it.},label={lst:x-blinky-out}]{out_blinky.txt}
